%% cuadro_ecuaciones.tex — Comparación de las 6 ecuaciones del DSGE (Galí-Monacelli,
%% forma reducida) contra su contraparte (explícita, parcial o inferida) en el
%% sistema de agentes LLM.
%%
%% Standalone: compilar este archivo directamente genera un PDF de solo esta tabla
%% (reutiliza preamble.tex de esta misma carpeta).

\documentclass[12pt, a4paper]{report}
%% preamble.tex — Paquetes y configuración de estilo académico

%% ─── Codificación y lenguaje ──────────────────────────────────────────────────
\usepackage[T1]{fontenc}
\usepackage[utf8]{inputenc}
\usepackage[spanish, es-nodecimaldot, es-noquoting, english]{babel}
\selectlanguage{spanish}

%% ─── Página y espaciado ───────────────────────────────────────────────────────
\usepackage[
  a4paper,
  top=3cm, bottom=3cm,
  left=3cm, right=2.5cm
]{geometry}
\usepackage[onehalfspacing]{setspace}

%% ─── Matemáticas ──────────────────────────────────────────────────────────────
\usepackage{amsmath, amssymb, amsthm}
\usepackage{mathtools}

%% Entornos matemáticos
\newtheorem{theorem}{Teorema}[chapter]
\newtheorem{proposition}[theorem]{Proposición}
\newtheorem{definition}[theorem]{Definición}
\newtheorem{remark}[theorem]{Observación}

%% ─── Figuras y tablas ─────────────────────────────────────────────────────────
\usepackage{graphicx}
\usepackage[export]{adjustbox}
\usepackage{booktabs}
\usepackage{tabularx}
\usepackage{longtable}
\usepackage{caption}
\usepackage{subcaption}
\usepackage{float}
\graphicspath{{figures/}}

%% ─── Código y listings ────────────────────────────────────────────────────────
\usepackage{listings}
\usepackage{xcolor}

\lstset{
  basicstyle=\ttfamily\small,
  breaklines=true,
  frame=single,
  captionpos=b,
  keywordstyle=\color{blue!70!black},
  commentstyle=\color{green!50!black},
  stringstyle=\color{orange!80!black}
}

%% ─── Bibliografía (biblatex + biber) ─────────────────────────────────────────
\usepackage[
  backend=biber,
  style=authoryear,
  natbib=true,
  sorting=nyt,
  maxbibnames=99,
  maxcitenames=2,
  uniquelist=false
]{biblatex}
\addbibresource{references.bib}

%% ─── Utilidades ───────────────────────────────────────────────────────────────
\usepackage{enumitem}
\usepackage{microtype}
\usepackage{soul}
\usepackage{csquotes}          % requerido por biblatex con babel

%% ─── Hipervínculos (SIEMPRE al final del preámbulo) ──────────────────────────
\usepackage[
  colorlinks=true,
  linkcolor=blue!70!black,
  citecolor=green!50!black,
  urlcolor=blue!80!black,
  pdftitle={Modelos DSGE vs Agentes LLM Heterogéneos},
  pdfauthor={},
  pdfsubject={Macroeconomía Computacional},
  pdfkeywords={DSGE, LLM, agentes, Argentina, simulación}
]{hyperref}

%% ─── Comandos custom para tesis en progreso ───────────────────────────────────
\newcommand{\placeholder}[1]{\textcolor{gray!70}{\textit{[#1]}}}
\newcommand{\todoc}[1]{\textcolor{red!80!black}{\textbf{[TODO: #1]}}}
\newcommand{\citeneeded}{\textsuperscript{\textcolor{red}{[cita]}}}

\begin{document}

\providecommand{\existeSi}[1]{\textcolor{green!45!black}{\textbf{#1}}}
\providecommand{\existeNo}[1]{\textcolor{red!70!black}{\textbf{#1}}}
\providecommand{\existeParcial}[1]{\textcolor{orange!80!black}{\textbf{#1}}}

\begingroup
\small
\setlength{\LTleft}{0pt}
\setlength{\LTright}{0pt}
\begin{longtable}{@{}c p{3.4cm} p{2.3cm} p{2.1cm} p{5.7cm}@{}}
\caption{Las seis ecuaciones del sistema DSGE (SOE-NK, forma reducida de
Galí-Monacelli 2005) y su contraparte —explícita, parcial o inferida— en el
sistema de agentes LLM.}
\label{tab:cuadro-ecuaciones-dsge-agentes} \\
\toprule
\textbf{\#} & \textbf{Ecuación DSGE} & \textbf{Origen real} & \textbf{¿Existe en agentes?} & \textbf{Cómo aparece en agentes} \\
\midrule
\endfirsthead

\multicolumn{5}{c}{\tablename\ \thetable\ -- continuación} \\
\toprule
\textbf{\#} & \textbf{Ecuación DSGE} & \textbf{Origen real} & \textbf{¿Existe en agentes?} & \textbf{Cómo aparece en agentes} \\
\midrule
\endhead

\midrule
\multicolumn{5}{r}{\textit{Continúa en la página siguiente}} \\
\endfoot

\bottomrule
\endlastfoot

1 &
\textbf{AR(1)} \newline
$r^{nat}_{t+1} = \rho_{rnat}\, r^{nat}_t + \varepsilon^{rnat}_t$ \newline
\textit{(demanda externa)}
&
Genérico NK, no específico de G\&M
&
\existeNo{No existe}
&
Ningún agente recibe ni genera una variable de ``tasa natural / demanda externa''. No hay contraparte.
\\
\addlinespace

2 &
\textbf{AR(1)} \newline
$s_{t+1} = \rho_s\, s_t + \varepsilon^{s}_t$ \newline
\textit{(cost-push)}
&
Genérico NK
&
\existeNo{No existe como proceso}
&
Las Firmas fijan \texttt{inflacion\_precio} con una heurística de \textit{pass-through} (\texttt{prompts.py:73,84}), pero no hay shock de costos explícito ni persistencia AR(1) --- es una decisión LLM período a período sin memoria estructural.
\\
\addlinespace

3 &
\textbf{AR(1)} \newline
$\phi_{t+1} = \rho_\phi\, \phi_t + \varepsilon^{\phi}_t$ \newline
\textit{(prima de riesgo / bimonetario)}
&
Adaptación propia (no G\&M)
&
\existeParcial{Existe (inferido)}
&
No hay shock exógeno. \texttt{brecha\_cambiaria} se actualiza como \texttt{prev}~$\times\,(1 + 0.10\cdot(\lambda_{mean}-0.30))$ (\texttt{simulation.py:172}). Es emergente: depende de cuánto se dolarizan los Hogares (\texttt{lambda\_usd}), decisión libre del LLM, no de un proceso estocástico calibrado.
\\
\addlinespace

4 &
\textbf{Regla de Taylor} \newline
$i_{t+1} = \rho_i\, i_t + (1-\rho_i)\left[\phi_\pi \pi_t + \phi_y y_t\right]$
&
Escudé / ARGEM, no G\&M
&
\existeSi{Sí, explícita (1:1)}
&
Se le da al Banco Central literalmente la misma fórmula con los mismos coeficientes ($\rho_i=0.80$, $\phi_\pi=1.50$, $\phi_y=0.125$) en su prompt (\texttt{prompts.py:49-53}). Es guía textual, no una restricción de código --- el LLM puede apartarse.
\\
\addlinespace

5 &
\textbf{Curva IS} \newline
$y_t = E_t y_{t+1} - \dfrac{1}{\sigma}(i_{t-1}-E_t\pi_{t+1}) + r^{nat}_t + \alpha\phi_t$
&
G\&M genuino (forma funcional)
&
\existeNo{No existe (inferido)}
&
\texttt{pbi\_gap} sale de \texttt{avg(produccion\_relativa)~-~1.0} (\texttt{simulation.py:169}). Las Firmas solo reciben la guía cualitativa ``$>1.0$ si demanda alta, $<1.0$ en recesión'' (\texttt{prompts.py:75,86}) --- sin tasa real, sin forward-looking, sin $\sigma$, sin canal de riqueza vía $\phi$. Es la brecha más grande entre ambos modelos.
\\
\addlinespace

6 &
\textbf{NKPC} \newline
$\pi_t = \beta E_t\pi_{t+1} + \kappa y_t + \alpha\phi_t + s_t$
&
G\&M genuino (forma funcional + $\kappa$ Calvo)
&
\existeParcial{Parcial}
&
Firmas usan reglas ad-hoc de \textit{pass-through}: tradable $\approx 0.7\,\Delta\%TC + 0.3\,\pi_{gral}$; no-tradable $\approx 0.9\,\pi_{gral} + 0.1\,\Delta\%TC$ (\texttt{prompts.py:72-73,83-84}). No hay término forward-looking $\beta E_t\pi_{t+1}$, no hay $\kappa$ (pendiente Calvo), no hay $y_t$ como insumo explícito.
\\

\end{longtable}
\endgroup

\end{document}
